% ---- pg-537 （书页 525） ----

在扑克牌的例子中, 将 16 张牌分为一组有一种方法. 有 32 767 种方法将 16 张牌分为两组(各组牌的张数可以不同). 有 7 141 686 种方法将牌分为三组(各组牌的张数可以不同), 等等.\footnote{将 $n$ 个对象分成 $k$ 个非空组的方法数即第二类的斯特林(Stirling)数, 其值为 $(1/k!)\sum_{j=0}^{k}(-1)^{k-j}\binom{k}{j}j^{n}$ (见文献[1]). 将这些数对 $k=1,2,\cdots,n$ 求和, 就得出将 $n$ 个对象分组的所有可能的方法数.} 显然, 由于时间的限制, 要从所有可能的分组中找出最优分组是不可能的. 即使大型计算机也难以处理这么多的计算量. 因此必须建立一些\mvbold{算法}, 利用它们可以找到好的, 但未必是最好的分组.

总之, 聚类分析的基本目标是发现项目(或变量)的自然分组方法. 为此, 我们必须先建立一个定量尺度, 借以量度对象之间的联系(相似性). 12.2 节用来讨论各种相似性量度. 在该节之后, 我们介绍几种比较通用的将对象分组的算法.

即使没有关于自然分组的确切概念, 我们也经常能凭眼睛将二维或三维图中的对象进行分组. 1.4 节中讨论过的星图和切尔诺夫脸可用于此种目的(见例 1.11 和 1.12). 还有一些方法能将高维观测值描绘为二维空间中的点, 因而从某种意义上说可使相似的对象彼此更接近. 这些方法将在 12.5$\sim$12.7 节中加以考虑.

\section{相似性量度}
\phantomsection\label{sec:12.2}

为了从复杂的数据集产生出比较简单的组结构, 多数做法都要求有一个“接近程度”或“相似性”的量度. 在选择相似性量度时, 通常会掺入相当大的主观性. 需要考虑的重要问题包括变量的性质(离散型、连续型、二值型), 测量值的尺度(名目尺度、顺序尺度、间隔尺度、比例尺度), 以及关于研究课题的知识.

当一些项目(单元或事例)被聚类时, 它们之间的接近程度通常用某种距离来表示. 另一方面, 变量通常根据相关系数或类似的联系量度来分组.

\mvsub{项目对的距离与相似性系数}

我们在第 1 章的 1.5 节曾讨论过距离概念. 根据\hyperref[eq:1-12]{式(1-12)}, 两个 $p$ 维观测值(项目) $\mathbf{x}=[x_1,\cdots,x_p]'$ 和 $\mathbf{y}=[y_1,\cdots,y_p]'$ 之间的欧氏距离(直线距离)为
\phantomsection\label{eq:12-1}
\begin{equation}
\begin{aligned}
d(\mathbf{x},\mathbf{y})&=\sqrt{(x_1-y_1)^2+(x_2-y_2)^2+\cdots+(x_p-y_p)^2}\\
&=\sqrt{(\mathbf{x}-\mathbf{y})'(\mathbf{x}-\mathbf{y})}
\end{aligned}
\tag{12-1}
\end{equation}
而同样两个观测值之间的统计距离[见\hyperref[eq:1-23]{式(1-23)}]为
\phantomsection\label{eq:12-2}
\begin{equation}
d(\mathbf{x},\mathbf{y})=\sqrt{(\mathbf{x}-\mathbf{y})'\mathbf{A}(\mathbf{x}-\mathbf{y})}
\tag{12-2}
\end{equation}
通常 $\mathbf{A}=\mathbf{S}^{-1}$, 其中 $\mathbf{S}$ 为样本方差和协方差矩阵. 但是, 没有关于不同组的先验知识, 这些样本量就无法计算. 由于这个原因, 在聚类问题中更倾向于采用欧氏距离.

另一种距离量度是闵可夫斯基(Minkowski)度量
\phantomsection\label{eq:12-3}
\begin{equation}
d(\mathbf{x},\mathbf{y})=\left[\sum_{i=1}^{p}|x_i-y_i|^{m}\right]^{1/m}
\tag{12-3}
\end{equation}
对 $m=1$, $d(\mathbf{x},\mathbf{y})$ 所量度的是 $p$ 维空间两个点之间的“城市街区”距离. 对 $m=2$, $d(\mathbf{x},\mathbf{y})$ 就是欧氏距离. 一般说来, 改变 $m$ 就改变了赋予较大或较小差值的权值.

两个额外的关于“距离”或相异性的常用量度为坎贝拉(Canberra)度量和捷卡诺夫斯基(Czekanowski)系数. 这两个量度都只对非负变量定义. 我们有
